\section{Discussion, Limitations, and Conclusion}
\label{sec:conclusion}

In this section, we discuss the generalization pathways and potential limitations of SABLE when extending beyond classification settings to generative models, followed by concluding remarks.

\subsection{Generalization to Generative Large Language Models}
A key technical dependency of SABLE's non-parametric routing is its reliance on pre-trained, task-specific classification weights (or centroids) $w_k$ to compute cosine similarities. In modern generative settings like Large Language Models (LLMs), there are typically no explicit task-specific classification heads; instead, models output tokens using a single, massive, shared vocabulary head. 

To generalize SABLE to generative LLM ensembling without violating our core "zero-parameter, zero-calibration data" philosophy, we propose utilizing a lightweight, frozen, pre-trained semantic text embedder (e.g., MiniLM) as a routing backbone. When a text prompt is received, the embedder projects the prompt into a low-dimensional semantic space. The similarity score $s_{k, b}$ is then computed by taking the cosine similarity of the prompt's embedding vector against pre-saved task-category centroid embeddings (derived from a small set of exemplar task instructions). Because sentence embedding is extremely fast, this preserves SABLE's stateless, real-time nature while extending its benefits to generative multi-expert serving pools. 

We explicitly acknowledge that our proposed generative LLM extension remains a conceptual, unproven blueprint and is a key limitation of the current work, as we do not provide empirical language modeling results here. To guide future empirical studies and make this pathway immediately actionable, we propose a precise evaluation setup: we will employ a LLaMA-3-8B base model adapted via LoRA adapters fine-tuned on four diverse text instruction tasks (e.g., SQL generation, creative writing, medical QA, and mathematical reasoning). For routing, we will use a frozen, off-the-shelf MiniLM encoder (22M parameters) to compute 384-dimensional prompt embedding vectors. Task centroids will be defined as the averaged embeddings of 16 representative task instructions (e.g., 'Write a SQL query...' or 'Solve the following math problem...'). Incoming user prompts will be projected on-the-fly and matched against these centroids via cosine similarity to compute SABLE blending coefficients. We will evaluate SABLE against standard weight-merging and MoE baselines, measuring both generation quality (ROUGE-L, accuracy) and serving latency on mixed prompt streams, verifying that SABLE maintains high generation fidelity under extreme prompt heterogeneity. This future validation will seek to empirically bridge the gap between classification tasks and full generative multi-expert serving pools, confirming SABLE's mathematical ensembling and efficiency advantages in the generative language domain.

\subsection{Physical Validation Scale and Multi-Layer Scaling Limitations}
Another limitation of this work is the scale of the physical convolutional neural network validation on real image datasets. In our physical CPU-based validation on MNIST and FashionMNIST, SABLE's activation blending is applied strictly at the final classification layer, leaving the convolutional base model static. While our 14-layer deep Analytical Coordinate Sandbox thoroughly validates SABLE's multi-layer ensembling across deep structures, the physical experiment is restricted to a single-layer projection setup. Consequently, the multi-layer physical validation on standard deep architectures (such as full ResNets, ViTs, or Transformers) remains an open area for future research. Indeed, to fully unlock standard-setting status and explore real-world high-dimensional computer vision tasks, we plan to evaluate SABLE on a pre-trained ViT-B/16 backbone adapted across 3-4 diverse tasks from the Visual Transfer Assessment Benchmark (VTAB) as part of our immediate future work. To make this future evaluation immediately actionable, we outline a concrete experimental blueprint: we will initialize a standard ViT-B/16 base model (86M parameters) pre-trained on ImageNet-21k, and fine-tune task-specific LoRA adapters ($r=8$, targeting query/value projection matrices) on four distinct VTAB categories (e.g., SVHN, CIFAR-100, DTD, and RESISC45). During test-time serving, a heterogeneous stream containing mixed images from these tasks will be routed dynamically using a support-split representation centroid in feature space, and SABLE Late Adaptation will blend the activation updates strictly at layers 10-12. This protocol will establish SABLE's scalability to real-world high-resolution, multi-layer vision applications. In deep architectures, applying SABLE across multiple hidden layers simultaneously could introduce representational drift or activation scale mismatches; while we analyze these challenges theoretically and propose mitigation techniques such as Late Adaptation and native normalization, physical verification of these multi-layer deep network configurations remains a critical next step for future empirical studies.

In future empirical scaling studies, practitioners can quantitatively track representational drift by measuring the cosine distance or centered kernel alignment (CKA) between the activations of SABLE and the unadapted base network across intermediate layers. This will provide a precise empirical map of drift trajectories, helping to optimize layer-dependent rank selection and routing depth bounds dynamically.

\subsection{Theoretical Limitations of Completely Zero-Data Centroids}
SABLE introduces a Completely Zero-Data Centroids construction method to enable dynamic ensembling under zero-calibration conditions. While empirically robust, this post-hoc parameter-based heuristic faces two key theoretical limitations:
\begin{enumerate}
    \item \textbf{Dual-Space Mismatch:} Classification weights are optimized parameters designed to project intermediate representations into logit space, whereas SABLE projects activation feature representations onto these weight vectors. Because weight spaces and feature activation spaces lie on fundamentally different manifolds and scales, measuring their cosine similarity directly constitutes a dual-space mismatch.
    \item \textbf{Vector Cancellation:} Discriminatively trained classification heads maximize class margins, often leading class-specific weights to point in orthogonal or negative directions. Averaging these class vectors row-wise to form a task-level centroid can trigger partial vector cancellation. This shrinks the centroid's norm and increases its sensitivity to high-frequency numerical noise.
\end{enumerate}
This mismatch explains the 5.80\% accuracy degradation of completely zero-data SABLE Soft compared to using a 16-sample support split (63.50\% vs 69.30\% at $r=10$). Future research can address these limitations by exploring weight-space normalization or projection-alignment layers to bridge the activation-parameter gap.

\paragraph{Generalizability of Weight L2-Normalization to Other Layers.} Our proposed weight-space L2-normalization trick works exceptionally well for final classification heads because they map representation features to class logits, where classes are structurally discriminative and symmetrical. An exciting question is whether this L2-normalization trick can be extended to intermediate projection matrices (e.g., self-attention query, key, value, or feed-forward projection layers) to construct zero-data task centroids for early-layer or mid-layer routing. For instance, in transformers, one could apply row-wise or column-wise L2-normalization to the raw low-rank update matrices $V_k = A_k B_k$ or individual matrices $A_k$ and $B_k$ to isolate the principal directional task coordinates, filtering out task-agnostic noise or magnitude offsets. This would enable zero-data centroid construction at any arbitrary layer depth, paving the way for multi-layer non-parametric routing without requiring any support-split activations, which is a promising direction for zero-data parameter-efficient serving.

\subsection{The Representational Blurring Paradox and Input-Space Routing Boundaries}
A core finding of our multi-layer physical MLP experiments (Section 4.4) is the +12.70\% performance surge of Single-Pass Early-Routing over standard 2-pass routing (65.20\% vs 52.50\%). This exposes the \textbf{Representational Blurring Paradox}: while a 2-pass configuration should theoretically perform better by using high-level penultimate representations, the unadapted layers of a jointly pre-trained base model compress and blur task-specific representations, degrading routing quality and injecting out-of-domain adapter noise.

For practitioners, this paradox highlights a critical deployability boundary of the SABLE framework:
\begin{itemize}
    \item \textbf{Input-Space Feasibility:} Single-Pass Early-Routing is highly effective when inputs are starkly separable in raw feature space (e.g., spatial image contours, or text prompts projected via high-fidelity frozen semantic embedders).
    \item \textbf{High-Dimensional Complex Noise:} When inputs are complex, high-dimensional, or noisy and lack clean input-level separability, Single-Pass Early-Routing suffers from extreme routing noise. In such scenarios, practitioners face a fundamental architectural trade-off: they must either accept the $2\times$ execution latency of 2-pass routing (while suffering from bottleneck representational blurring), introduce an external routing backbone (violating the pure single-model paradigm), or restrict adaptation exclusively to late-stage layers via SABLE's default Late Adaptation (Mid-Layer Routing) configuration.
\end{itemize}
Highlighting this boundary ensures that SABLE can be deployed with realistic, scientifically grounded design expectations across diverse production settings.

\subsection{Storage Scalability and PEFT Architecture Extensions}
A critical consideration for SABLE's deployment is the storage and memory bandwidth scalability of maintaining a pool of $K$ specialized low-rank experts. Although Top-$M$ expert pruning successfully bounds inference FLOPs and active execution complexity to $O(M)$ by executing only the most relevant experts, the system must still store the full pool of $K$ adapters. In massive multi-task settings where $K$ scales to dozens or hundreds of tasks, downloading and retaining these adapters in memory could eventually introduce a storage and bandwidth bottleneck. SABLE's PEFT parameterization heavily mitigates this compared to full-parameter ensembling: since each $r=8$ adapter represents $<1\%$ of the base model's capacity (e.g., only 1.2\% parameter overhead per expert in our sandbox), storing 50 SABLE experts requires less than 50\% of a single extra model's storage, whereas storing 50 full-parameter experts is completely prohibitive. For extreme scalability, future work can explore dynamically streaming adapter weights from host RAM to GPU VRAM on-the-fly, or using parameter-sharing layers across experts. Indeed, modern multi-tenant PEFT serving engines (such as Punica, S-LoRA, or vLLM) provide highly optimized, industrial-grade solutions to handle distributed adapter storage and dynamic loading. These systems manage a massive pool of inactive adapters in cheap host CPU memory or disk, and dynamically prefetch active adapters into GPU VRAM using tiered memory hierarchies and page-aligned adapter cache managers (e.g., TokenAttention or PagedAttention structures). By overlapping GPU execution of the current batch with the asynchronous, PCIe-bound streaming of next-batch adapters, these serving frameworks completely eliminate dynamic loading overhead. SABLE's stateless architecture is naturally designed to integrate with these multi-tenant engines, allowing practitioners to store hundreds of adapters with negligible memory footprint and run stateless ensembling with zero systems overhead.

Furthermore, SABLE's activation blending algebra can be naturally extended to alternative Parameter-Efficient Fine-Tuning (PEFT) architectures beyond low-rank LoRA matrices. For instance, in $(IA)^3$ architectures where task-specific updates are parameterized as element-wise scaling vectors $l_k \in \mathbb{R}^D$ applied to key, value, and feed-forward activations, SABLE's activation ensembling can be formulated as a dynamic, sample-wise convex combination of scaling vectors: $Y_b = X_b \cdot \sum_k \alpha_{k, b} l_k$. Similarly, for Prefix Tuning or Prompt Tuning, the key and value prefix matrices or prompt embeddings can be dynamically interpolated in the forward pass using SABLE routing coefficients. This demonstrates SABLE's mathematical flexibility and paves the way for a unified activation-space ensembling paradigm across diverse PEFT methods.

\subsection{Mitigating Capacity Bottlenecks via Layer-Dependent Hybrid-Rank Selection}
Our physical CNN and ResNet-18 foundation experiments (Sections 4.3 and 4.4) expose a capacity bottleneck when rank $r$ is constrained in low-dimensional output projection layers, with SABLE's performance dropping significantly at $r \le 4$. This highlights a key sensitivity where low-rank constraints can be overly restrictive in layers with low intrinsic dimensionality.

To mitigate this capacity bottleneck while maintaining global parameter efficiency, we propose a \textbf{Layer-Dependent Hybrid-Rank Selection Protocol} as a core design recommendation:
\begin{itemize}
    \item \textbf{Low-Rank Hidden Layers:} Keep attention projection and hidden layer adapters at strict, aggressive low rank (e.g., $r=4$ or $r=8$) to minimize intermediate activation memory and parameter overhead.
    \item \textbf{Full-Rank Output Projection Blocks:} Configure final classification heads or vocabulary output projection layers at full task rank (e.g., $r = K$, where $K$ is the number of active task experts). This preserves representation capacity where logit discriminability is resolved, preventing capacity choke-points with minimal impact on total parameter storage.
\end{itemize}
Formalizing this hybrid rank protocol optimizes SABLE's parameter-efficiency/representation trade-off, ensuring high task fidelity across all projection boundaries.

\subsection{Concluding Remarks}
In this work, we introduced \textbf{SABLE (Sample-wise Activation Blending of Low-Rank Experts)}, a minimalist, stateless, and network-level framework for test-time dynamic model merging. Driven by the principle of Occam's razor, SABLE shifts the ensembling step from parameter space to activation space, allowing mathematically exact ensembling on a per-sample basis during the forward pass. 

SABLE completely strips away the stateful dynamic buffers, sorting steps, and scheduling latency associated with current systems-level solutions like Micro-Batch Homogenization (MBH). Through extensive evaluations in the Analytical Coordinate Sandbox, we demonstrated that SABLE is completely immune to stream heterogeneity. Specifically, SABLE configured with our default Late Adaptation (Mid-Layer Routing) achieves a flatline joint mean accuracy of \textbf{68.10\%} across both homogeneous and heterogeneous streams, natively outperforming the state-of-the-art, systems-heavy PFSR+MBH pipeline (\textbf{67.20\%}) in a single, stateless forward pass. Even our simpler Early-Layer Routing configuration delivers a stable joint mean accuracy of \textbf{66.60\%}, matching the performance of the systems-heavy MBH pipeline within a negligible \textbf{0.60\%}.

Our work demonstrates that when faced with stream-induced degradation, we do not need to build complex systems wrappers around deep networks. Instead, a simple change in ensembling algebra combined with parameter-efficient representations can resolve the problem natively, returning model serving to its clean, stateless, and reproducible roots.

For future work, we plan to implement our generative LLM routing extension and evaluate SABLE on massive multi-task language benchmarks, moving even closer to the absolute lower-bound of architectural complexity in foundation model ensembling.
